% from ega0-22-fr.tex
\oldpage[0\textsubscript{IV-1}]{206}

  soit $\mathfrak L$ un idéal de définition de $C$~; alors les conditions a)
  et b) de (i) équivalent aussi à~:
  \begin{enumerate}
    \item[c)] L'homomorphisme de $(C/\mathfrak L)$-modules
    \[
      (\mathfrak J/\mathfrak J^2)\otimes_C(C/\mathfrak L)
      \longrightarrow\Omega^1_{B/A}\otimes_B(C/\mathfrak L)
    \]
    est inversible à gauche.
  \end{enumerate}
\end{enumerate}
\end{corollary}

\begin{proof}
Les hypothèses entraînent que sur $\Omega^1_{B/A}$ la topologie est déduite de
celle de $B$ (20.4.5), donc (i) résulte de (22.6.1) et de (19.1.7). D'autre
part, rappelons (20.4.9) que puisque $B$ est une $A$-algèbre formellement
lisse, $\Omega^1_{B/A}$ est un $B$-module formellement projectif, donc
$\Omega^1_{B/A}\otimes_BB_\lambda$ est un $B_\lambda$-module projectif pour
tout $\lambda$ (19.2.4)~; l'équivalence de c) avec a) et b) lorsque $C$ est
préadmissible résulte alors de (19.1.9).
\end{proof}

En particulier~:

\begin{corollary}[22.6.3]
\label{0.22.6.3-fr}
Soient $A$ un anneau, $B$ une $A$-algèbre, $\mathfrak J$ un idéal de $B$,
$C=B/\mathfrak J$~; on suppose $A$, $B$, $C$ munis des topologies discrètes
et que $B$ est une $A$-algèbre formellement lisse. Pour que $C$ soit une
$A$-algèbre formellement lisse, il faut et il suffit que l'homomorphisme
canonique (22.6.1.1) soit inversible à gauche.
\end{corollary}

On notera que puisque toute $A$-algèbre $C$ est quotient d'une algèbre de
polynômes $A[T_\alpha]_{\alpha\in I}$, et que cette dernière est formellement
lisse (19.3.3), le critère (22.6.3) permet en principe de reconnaître si $C$
est une $A$-algèbre formellement lisse.

\begin{proposition}[22.6.4]
\label{0.22.6.4-fr}
Soient $A$ un anneau, $B$ une $A$-algèbre formellement lisse (pour les
topologies discrètes), $\mathfrak J$ un idéal de $B$, $C=B/\mathfrak J$~; on
suppose que $\mathfrak J/\mathfrak J^2$ est un $C$-module de type fini. Soient
$\mathfrak p$ un idéal premier de $C$, $k(\mathfrak p)$ le corps résiduel de
$C_{\mathfrak p}$, $\mathfrak p'$ l'image réciproque de $\mathfrak p$ dans
$B$, $\mathfrak q$ l'idéal premier de $A$ image réciproque de $\mathfrak p'$.
Les conditions suivantes sont équivalentes~:
\begin{enumerate}
  \item[a)] $C_{\mathfrak p}$ est une $A_{\mathfrak q}$-algèbre formellement
  lisse (pour les topologies discrètes).
  \item[a')] $C_{\mathfrak p}$ est une $A$-algèbre formellement lisse (pour
  les topologies discrètes).
  \item[b)] L'homomorphisme canonique
  \begin{equation}
  \label{0.22.6.4.1-fr}
    (\mathfrak J/\mathfrak J^2)\otimes_Ck(\mathfrak p)
    \longrightarrow\Omega^1_{B/A}\otimes_Bk(\mathfrak p)
    \tag{22.6.4.1}
  \end{equation}
  est injectif.
  \item[c)] Il existe $f\in C-\mathfrak p$ tel que $C_f$ soit une
  $A$-algèbre formellement lisse.
\end{enumerate}
Si de plus $B$ est noethérien, les conditions précédentes sont encore
équivalentes à~:
\begin{enumerate}
  \item[d)] $C_{\mathfrak p}$ est une $A_{\mathfrak q}$-algèbre formellement
  lisse pour la topologie $\mathfrak p$-préadique sur $C_{\mathfrak p}$ et la
  topologie discrète (ou la topologie $\mathfrak q$-préadique) sur
  $A_{\mathfrak q}$.
\end{enumerate}
\end{proposition}

\begin{proof}
Comme $A_{\mathfrak q}$ est une $A$-algèbre formellement lisse (19.3.5, (iv)
et (ii)), on sait déjà que a) et a') sont équivalentes (19.3.5, (ii)). On a
$C_{\mathfrak p}=B_{\mathfrak p'}/\mathfrak J B_{\mathfrak p'}$, et
$B_{\mathfrak p'}$ est une $A$-algèbre formellement lisse pour la topologie
discrète (19.3.5, (iv))~; en outre
$\Omega^1_{B_{\mathfrak p'}/A}=(\Omega^1_{B/A})_{\mathfrak p'}$ (20.5.9).
L'équivalence de a') et b) résulte alors de l'application de (22.6.3) à
$B_{\mathfrak p'}$ et $C_{\mathfrak p}$, en tenant compte de ce que
$\Omega^1_{B/A}$ est un $B$-module projectif (20.4.9) et (19.2.11), et que
$\mathfrak J/\mathfrak J^2$ est un $B$-module de type fini, et en utilisant
(19.1.12). L'application de (19.1.12) prouve aussi l'équivalence de b) et c).
Enfin, notons que puisque $B_{\mathfrak p'}$ est une
$A_{\mathfrak q}$-algèbre formellement lisse pour les topologies

